% !TeX root = ../../../main.tex
\subsubsection{Data and Extraction Pipeline}
We developed a four-stage pipeline (Figure \ref{ch06:fig:flow}) for case report detection, Sepsis-3 phenotyping, clinical finding extraction, and timestamp determination, applied to two data sources: (1) the absolute timeline annotations from \cite{ch06:frattallone2024using}, and (2) the PubMed Open Access subset \citep{pmc2024open} (see Table \ref{ch06:tab:dataset-summary}).


% !TeX root = ../../../main.tex
\begin{table}[htbp]
\caption{Dataset summary. Bold cells indicate contributed labeling}\ChSixFloatLabel{ch06:tab:dataset-summary}
\ChSixTableBodyStart
\begin{ChSixFloatBody}




\centering
% \begin{tabular}{|l|l|l|l|}
\begin{adjustbox}{width=0.8\textwidth}
\begin{tabular}{|>{\raggedright\arraybackslash}p{3.5cm}|>{\raggedright\arraybackslash}p{3.5cm}|>{\raggedright\arraybackslash}p{3.5cm}|>{\raggedright\arraybackslash}p{3.5cm}|}
\hline
\textbf{Corpus}                  & \centering \textbf{i2m4}             & \centering \textbf{\texttt{sepsis-40}}& \textbf{\qquad \texttt{sepsis-100}}       \\
\hline
\textbf{N (unlabeled)}          & \multicolumn{1}{r|}{310;  331,794}             & \multicolumn{1}{r|}{2,139}                     & \multicolumn{1}{r|}{2,139}                      \\ \hline
\textbf{N (labeled)}            & \multicolumn{1}{r|}{20}                         & \multicolumn{1}{r|}{40}                       & \multicolumn{1}{r|}{100}                       \\ \hline
\textbf{Source}                 & i2b2; MIMIC IV           & PMOA       & PMOA       \\ \hline
\textbf{Note type}              & discharge summary        & case report              & case report               \\ \hline
\textbf{Labeler}                & \textbf{clinicians}               & \textbf{clinicians} (Gold)& \textbf{GPT-5} (Bronze)\\ \hline
\textbf{Phenotyper}             & \textbf{clinicians}      & \textbf{clinicians}      & \textbf{QwQ-32B-Preview}             \\ \hline
\textbf{Methods}                &                & GPT-5& GPT-5\\ 
                                & DeepSeek-R1& DeepSeek-R1& DeepSeek-R1\\ 
                                & L3.3 70B Instruct                  & L3.3 70B Instruct                & L3.3 70B Instruct                 \\ 
                                & L3.1 70B Instruct                 & L3.1 70B Instruct                & L3.1 70B Instruct   \\ 
                                & Mixtral 8x7B & Mixtral 8x7B                & Mixtral 8x7B                   \\ 
                                & \quad Instruct v0.1 & \quad Instruct v0.1                & \quad Instruct v0.1                   \\ 
                                & L3.1 8B Instruct             & L3.1 8B Instruct                & L3.1 8B Instruct                  \\ \hline
\parbox{3cm}{\textbf{Ablations}} & No role                  & No role                  & No role                   \\ 
Base: & No conjunction & No conjunction & No conjunction
                                \\ 
L3.3 70B Instruct                        & \quad instruction &  \quad instruction & \quad instruction \\
                                 % &&&  \\
                                
                                & No few-shot prompting    & No few-shot prompting    & No few-shot prompting     \\ \hline
\textbf{Sensitivity analyses}   & With upper bound         & With upper bound         & With upper bound          \\ 
                                & With upper bound + 
                                                         & With upper bound +  
                                                                                      & With upper bound + \\
                                & \quad i2b2 event typing & \quad i2b2 event typing & \quad i2b2 event typing \\
                                                                                      \hline
\end{tabular}
\end{adjustbox}

\end{ChSixFloatBody}
\end{table}


\paragraph{Datasets}
\label{ch06:section:datasets}

\ChSixParagraph{i2m4.} This dataset contains 20 discharge summaries (15 i2b2, 5 MIMIC-IV) with physician-annotated temporal concepts from \cite{ch06:frattallone2024using}. We used interval lower bounds as event times and only probabilistic annotations. For Sepsis-3 phenotyping, we used a prompt describing qSOFA criteria: suspected/confirmed infection with blood pressure/respiratory rate/mental status abnormalities, requesting binary output (1 for sepsis, 0 otherwise).
To develop a computational approach, we used a prompt describing qSOFA criteria (Appendix \ref{ch06:apd:phenotype-prompt}) with QwQ-32B-Preview achieving 100\% accuracy and Llama-3.1-8B-Instruct achieving 65\% accuracy.

\ChSixParagraph{PMOA (\texttt{sepsis-40}, \texttt{sepsis-10}, \texttt{sepsis-100}).}
From 1.5M publications (version April 25, 2024), we identified case reports via regex (\texttt{case (report|presenta)} and \texttt{year-?old}), then filtered for sepsis using \texttt{(sepsi|septic)} and \texttt{(critical|intensive) care}.
GPT-4o and Llama-3.1-8B screened for Sepsis-3, yielding 2,139 cases.
We manually annotated 40 PMOA case reports for the gold-standard \texttt{sepsis-40} evaluation set, yielding $3,229$ clinician-annotated reference events.
The \texttt{sepsis-10} subset is a 10-report subset of \texttt{sepsis-40} used for clinician annotation quality review, hallucination audit, multi-pass self-consistency, and sepsis onset analysis.
We also sampled 100 non-overlapping PMOA case reports for \texttt{sepsis-100}, which is used as a bronze-standard proxy-reference set annotated by GPT-5 to evaluate extraction behavior at larger scale rather than as clinical ground truth.
The full case identifiers for \texttt{sepsis-40} and \texttt{sepsis-10} are provided in Appendix~\ref{ch06:apd:subset_ids}.
Figure~\ref{ch06:fig:flow} (right) shows phenotyping confusion matrices.

\textbf{Phenotyping output parsing and clinician audit.}
In practice, some LLMs produce verbose rationales rather than a minimal binary output. We therefore parsed phenotyping decisions using a rule-based extraction step that searches model responses for explicit affirmative/negative indicators of meeting sepsis criteria and maps them to a binary label. If primary extraction did not yield a valid label, a fallback strategy examined the final line of the response, and discordant parses were flagged for manual inspection. We additionally performed clinician audit on sampled and borderline cases to verify alignment with Sepsis-3 intent under information available in narratives.


\paragraph{Clinical Finding Extraction and Temporal Assignment}
We queried six LLMs (GPT-5, DeepSeek-R1-UD-IQ1\_S, Llama-3.3-70B-Instruct, Llama-3.1-70B-Instruct, Mixtral-8x7B-Instruct v0.1, Llama-3.1-8B-Instruct) using a structured prompt with few-shot examples to extract event-time pairs. The prompt instructs models to: assign admission time as $t=0$, use negative times for pre-admission events, split conjunctions, and capture event start times for durations (full prompt in Appendix \ref{ch06:apd:prompt}). For sensitivity analysis, we tested ablations (no-role, zero-shot, no-expansion) and augmentations (interval bounds, i2b2 typing).

For i2m4, we compared against a BERT baseline adapted from \cite{ch06:frattallone2024using}, using a 3-layer MLP regressor with L1 loss for time prediction. We filtered events to $\pm$2 weeks of admission for training stability (3,472$\rightarrow$3,350 findings) and used 5-fold cross-validation.


